\documentclass[12pt,a4paper]{article}

% Drake_preamble.tex - LaTeX Preamble Template
% 作者: 謝慕揚, MD, PhD, FESC
% 
% 使用說明:
% 1. 表格請使用 longtable，不要有垂直分隔線 |
% 2. 表格內換行使用 \newline，不要使用 \\
% 3. 藥物名稱、檢驗、檢查請維持英文
% 4. Reference 格式請使用 \href 加上驗證過的 DOI

% Including essential packages for Chinese typesetting
\usepackage{xeCJK}
\usepackage{fontspec}
% Configuring Chinese font with Noto Serif CJK TC, fallback to SimSun
\IfFontExistsTF{Noto Serif CJK TC}{
  \setCJKmainfont{Noto Serif CJK TC}
  \setCJKsansfont{Noto Serif CJK TC}
  \setCJKmonofont{Noto Serif CJK TC}
}{\setCJKmainfont{SimSun}
  \setCJKsansfont{SimSun}
  \setCJKmonofont{SimSun}
}

% Including other necessary packages
\usepackage{geometry}
\usepackage{titlesec}
\usepackage{enumitem}
\usepackage{xcolor}
\usepackage{colortbl}
\usepackage{tcolorbox}
\usepackage{booktabs}
\usepackage{longtable}
\usepackage{array}
\usepackage{graphicx}
\usepackage{tikz}
\usepackage{fancyhdr}
\usepackage{multicol}
\usepackage{datetime2}
\usepackage{hyperref}
\usepackage{mdframed}
\usepackage{amsmath}
\usepackage{amssymb}
\usepackage{multirow}

\hypersetup{
    colorlinks=true,
    linkcolor=emergency_blue,
    citecolor=emergency_blue,
    urlcolor=emergency_blue,
    pdfborder={0 0 0}
}

% Defining page geometry
\geometry{
  top=2.5cm,
  bottom=2.5cm,
  left=2cm,
  right=2cm,
  headheight=25pt
}

% Adjusting topmargin to compensate for increased headheight
\addtolength{\topmargin}{-10pt}

% Defining colors
\definecolor{emergency_red}{RGB}{186,24,27}
\definecolor{emergency_blue}{RGB}{0,114,188}
\definecolor{emergency_gray}{RGB}{120,120,120}
\definecolor{highlight_yellow}{RGB}{255,250,205}

% Setting title formats
\titleformat{\section}[block]{\Large\bfseries\color{emergency_red}}{\thesection}{10pt}{}
\titleformat{\subsection}[block]{\large\bfseries\color{emergency_blue}}{\thesubsection}{8pt}{}
\titleformat{\subsubsection}[block]{\normalsize\bfseries\color{emergency_gray}}{\thesubsubsection}{6pt}{}

% Defining custom environment for important notes
\newmdenv[
    linecolor=emergency_red,
    backgroundcolor=highlight_yellow,
    linewidth=2pt,
    topline=true,
    bottomline=true,
    leftline=true,
    rightline=true,
    innertopmargin=10pt,
    innerbottommargin=10pt,
    innerleftmargin=10pt,
    innerrightmargin=10pt
]{importantbox}

% Defining note command with tcolorbox
\newcommand{\note}[1]{
    \par
    \noindent
    \begin{tcolorbox}[
        colback=emergency_blue!5!white,
        colframe=emergency_blue,
        title=\textbf{重點提醒},
        fonttitle=\bfseries,
        boxrule=0.5pt,
        arc=3pt,
        left=6pt,
        right=6pt,
        top=6pt,
        bottom=6pt
    ]
    #1
    \end{tcolorbox}
}

% Key findings box
\newcommand{\keyfinding}[1]{
    \par
    \noindent
    \begin{tcolorbox}[
        colback=yellow!10!white,
        colframe=orange!75!black,
        title=\textbf{核心發現摘要},
        fonttitle=\bfseries,
        boxrule=0.5pt,
        arc=3pt
    ]
    #1
    \end{tcolorbox}
}

% Defining custom date format for ISO 8601
\DTMnewdatestyle{iso}{
  \renewcommand{\DTMdisplaydate}[4]{\number##1-\number##2-\number##3}
  \renewcommand{\DTMDisplaydate}[4]{\number##1-\number##2-\number##3}
}
\DTMsetdatestyle{iso}
\newcommand{\mydate}{\DTMtoday}


% Custom header for this document
\pagestyle{fancy}
\fancyhf{}
\fancyhead[L]{\textcolor{emergency_red}{NEJM AI 教學講義}}
\fancyhead[R]{\textcolor{emergency_blue}{AI in Medical Education}}
\fancyfoot[C]{\textcolor{emergency_gray}{\thepage}}
\renewcommand{\headrulewidth}{1pt}
\renewcommand{\headrule}{\hbox to\headwidth{\color{emergency_red}\leaders\hrule height \headrulewidth\hfill}}

\begin{document}

%============================================================
% Title Page
%============================================================
\begin{center}
{\LARGE\bfseries\textcolor{emergency_red}{NEJM AI}}\\[0.3cm]
{\Large\textcolor{emergency_blue}{Case Study}}\\[0.5cm]
{\Large\textbf{Assessment of Short-Answer Questions by ChatGPT in a Medical School Course}}\\[0.3cm]
{\large 使用 ChatGPT 評量醫學院課程簡答題}\\[0.5cm]
{\normalsize \href{https://doi.org/10.1056/AIcs2500239}{NEJM AI 2026;3(2). DOI: 10.1056/AIcs2500239}}\\[0.3cm]
{\normalsize 作者：Kuling G, Pullman S, Vasilev D, et al. (Harvard Medical School)}\\[0.3cm]
{\normalsize 整理：謝慕揚 MD, PhD, FESC \quad 日期：\mydate}
\end{center}

\vspace{0.5cm}

%============================================================
\section{研究背景 (Background)}
%============================================================

\subsection{教育學原理}

\begin{itemize}[leftmargin=2em]
    \item \textbf{頻繁的低風險測驗 (Frequent low-stakes testing)} 是促進學習的有力工具
    \item \textbf{簡答題 (Short-answer questions)} 能促進批判性思考 (critical thinking)
    \item 然而，簡答題在形成性評量 (formative testing) 中\textbf{使用不足}
    \item 原因：需要\textbf{耗時的人工評分}
\end{itemize}

\subsection{研究目的}

\begin{importantbox}
\textbf{主要目標：}比較 AI (GPT-4o) 與人類評分者對醫學生開放式簡答題的評分一致性。

\textbf{核心問題：}AI 能否減輕簡答題評分的負擔，同時維持評分品質？
\end{importantbox}

%============================================================
\section{研究方法 (Methods)}
%============================================================

\subsection{研究設計}

\begin{longtable}{p{4cm}p{10cm}}
\toprule
\textbf{項目} & \textbf{說明} \\
\midrule
\endfirsthead
研究類型 & Retrospective case study \\
\midrule
AI 模型 & GPT-4o (Generative Pretrained Transformer 4o) \\
\midrule
研究對象 & 169 位 Harvard Medical School 一年級醫學生 \\
\midrule
課程 & 2023-2024 學年 Foundations course \\
\midrule
人類評分者 & 10 位受過訓練的評分者 \\
\bottomrule
\end{longtable}

\subsection{評分標準}

學生答案依據兩個獨立標準進行評分：

\begin{enumerate}[leftmargin=2em]
    \item \textbf{Factual accuracy（事實正確性）}：答案內容是否正確
    \item \textbf{Completeness（完整性）}：答案是否涵蓋所有重要面向
\end{enumerate}

\subsection{資料集設計}

\begin{longtable}{p{4cm}p{5cm}p{5cm}}
\toprule
\textbf{資料集} & \textbf{Multi-grader Set} & \textbf{Single-grader Set} \\
\midrule
\endfirsthead
樣本數 & n = 822 responses & n = 8,964 responses \\
\midrule
評分方式 & 10 位人類評分者獨立評分\textbf{相同}答案 & 每個答案由\textbf{一位}人類評分者評分 \\
\midrule
用途 & Prompt 優化與教學原則校準 & 模擬真實世界條件進行模型評估 \\
\bottomrule
\end{longtable}

\subsection{Prompt Engineering 策略}

\note{
GPT prompts 經過\textbf{反覆優化 (iteratively refined)}，以達成：
\begin{itemize}[leftmargin=2em]
    \item \textbf{Semantic fairness（語意公平性）}：避免因用詞不同而影響評分
    \item \textbf{Pedagogical alignment（教學原則一致性）}：符合教育學評量標準
    \item \textbf{Linguistic flexibility（語言彈性）}：容許學生不同的表達方式
\end{itemize}
}

\subsection{統計分析}

\textbf{Quadratic-weighted Cohen's Kappa：}
\begin{itemize}[leftmargin=2em]
    \item 用於衡量 AI 與人類評分的一致性程度
    \item 考慮評分差異的大小（差異越大，懲罰越重）
    \item Kappa 值解讀：
    \begin{itemize}
        \item < 0.20：Poor agreement
        \item 0.21-0.40：Fair agreement
        \item 0.41-0.60：Moderate agreement
        \item 0.61-0.80：Substantial agreement
        \item 0.81-1.00：Almost perfect agreement
    \end{itemize}
\end{itemize}

%============================================================
\section{研究結果 (Results)}
%============================================================

\subsection{Multi-grader Dataset 結果}

\begin{center}
\begin{tabular}{p{5cm}p{4cm}p{4cm}}
\toprule
\textbf{評分標準} & \textbf{Quadratic-weighted Kappa} & \textbf{一致性程度} \\
\midrule
Factual accuracy & $0.443 \pm 0.127$ & Moderate \\
\midrule
Completeness & $0.429 \pm 0.145$ & Moderate \\
\bottomrule
\end{tabular}
\end{center}

\keyfinding{
\textbf{重要發現：}94\% 的 GPT 生成的 factual accuracy 評分落在人類評分者評分範圍內。

這表示 GPT 的評分變異性與人類評分者之間的變異性相當。
}

\subsection{Single-grader Dataset 結果}

\begin{center}
\begin{tabular}{p{5cm}p{4cm}p{4cm}}
\toprule
\textbf{評分標準} & \textbf{Adjusted Kappa}$^*$ & \textbf{一致性程度} \\
\midrule
Factual accuracy & 0.741 & Substantial \\
\midrule
Completeness & 0.655 & Substantial \\
\bottomrule
\end{tabular}
\end{center}

{\small $^*$Adjusted for rater variability using an attenuation factor}

\begin{importantbox}
\textbf{臨床意義：}超過 \textbf{85\%} 的 GPT 評分與人類評分者的分數差距在 \textbf{1 分以內}。

在實際應用情境中，GPT-4o 與人類評分者達到 \textbf{substantial agreement}。
\end{importantbox}

\subsection{結果視覺化比較}

\begin{center}
\begin{tabular}{p{4cm}p{4cm}p{4cm}}
\toprule
 & \textbf{Factual Accuracy} & \textbf{Completeness} \\
\midrule
Multi-grader Kappa & 0.443 (Moderate) & 0.429 (Moderate) \\
\midrule
Single-grader Kappa & 0.741 (Substantial) & 0.655 (Substantial) \\
\midrule
GPT 分數在人類範圍內 & 94\% & -- \\
\midrule
GPT 與人類差距 $\leq$1 分 & >85\% & >85\% \\
\bottomrule
\end{tabular}
\end{center}

%============================================================
\section{討論 (Discussion)}
%============================================================

\subsection{為何 Single-grader Kappa 較高？}

\begin{itemize}[leftmargin=2em]
    \item Multi-grader dataset 反映了\textbf{人類評分者之間的變異性}
    \item 即使是受過訓練的人類評分者，對同一答案的評分也存在差異
    \item Single-grader dataset 使用 \textbf{attenuation factor} 校正評分者變異性
    \item 校正後的 kappa 值更能反映 GPT 與「理想人類評分」的一致性
\end{itemize}

\subsection{Prompt Engineering 的重要性}

本研究強調了 \textbf{iterative prompt engineering} 的關鍵角色：

\begin{longtable}{p{4cm}p{10cm}}
\toprule
\textbf{原則} & \textbf{實踐方式} \\
\midrule
\endfirsthead
Pedagogical principles & 將教育學評量原則融入 prompt 設計 \\
\midrule
Semantic fairness & 確保不同用詞但意義相同的答案獲得相同評分 \\
\midrule
Linguistic flexibility & 容許學生使用不同的專業術語或表達方式 \\
\midrule
Iterative refinement & 根據 multi-grader dataset 的回饋不斷優化 prompt \\
\bottomrule
\end{longtable}

\subsection{研究限制}

\begin{itemize}[leftmargin=2em]
    \item 單一機構研究（Harvard Medical School）
    \item 僅評估一門基礎課程
    \item 未評估長篇申論題或臨床推理題
    \item AI 評分的長期教育效果尚待研究
\end{itemize}

%============================================================
\section{臨床與教育應用}
%============================================================

\subsection{潛在應用場景}

\begin{enumerate}[leftmargin=2em]
    \item \textbf{形成性評量 (Formative Assessment)}
    \begin{itemize}
        \item 提供即時回饋給學生
        \item 增加低風險測驗的頻率
        \item 減輕教師評分負擔
    \end{itemize}

    \item \textbf{大規模開放式線上課程 (MOOCs)}
    \begin{itemize}
        \item 使簡答題在大規模教育中變得可行
        \item 提供個人化學習回饋
    \end{itemize}

    \item \textbf{自我評量工具}
    \begin{itemize}
        \item 學生可使用 AI 進行自我練習
        \item 獲得即時的正確性與完整性回饋
    \end{itemize}
\end{enumerate}

\subsection{實施建議}

\note{
\textbf{使用 AI 評分系統的建議：}
\begin{enumerate}[leftmargin=2em]
    \item 進行機構特定的 prompt 優化與驗證
    \item 建立人類監督機制（特別是高風險評量）
    \item 定期校準 AI 與人類評分的一致性
    \item 向學生說明 AI 評分的角色與限制
    \item 保留人類複審的管道
\end{enumerate}
}

%============================================================
\section{Key Points 重點整理}
%============================================================

\begin{importantbox}
\begin{enumerate}[leftmargin=2em]
    \item \textbf{研究目的：}評估 GPT-4o 評分醫學生簡答題與人類評分者的一致性

    \item \textbf{評分標準：}Factual accuracy（事實正確性）與 Completeness（完整性）

    \item \textbf{主要發現：}
    \begin{itemize}
        \item Multi-grader dataset：Kappa 約 0.43-0.44（moderate agreement）
        \item Single-grader dataset：Kappa 0.66-0.74（substantial agreement）
        \item 94\% GPT 評分落在人類評分範圍內
        \item >85\% GPT 評分與人類差距 $\leq$1 分
    \end{itemize}

    \item \textbf{成功關鍵：}Iterative prompt engineering 結合教學原則與語言彈性

    \item \textbf{應用價值：}AI 評分系統有潛力減輕簡答題評分負擔，使此類題型更廣泛應用於醫學教育

    \item \textbf{注意事項：}目前達到 moderate-to-substantial agreement，高風險評量仍需人類監督
\end{enumerate}
\end{importantbox}

%============================================================
\section{延伸思考}
%============================================================

\subsection{AI 在醫學教育的未來角色}

\begin{itemize}[leftmargin=2em]
    \item \textbf{個人化學習路徑：}根據學生答題表現，AI 可識別知識缺口並推薦學習資源
    \item \textbf{即時回饋系統：}學生可在練習時獲得立即的評分與解釋
    \item \textbf{教師時間重新分配：}減少機械性評分工作，讓教師專注於高價值教學互動
    \item \textbf{評量多樣化：}使過去因評分困難而避免的題型重新成為選項
\end{itemize}

\subsection{倫理考量}

\begin{itemize}[leftmargin=2em]
    \item \textbf{透明度：}學生是否應被告知其答案由 AI 評分？
    \item \textbf{公平性：}AI 評分是否對特定背景學生存在偏見？
    \item \textbf{隱私：}學生答案資料如何被處理與保護？
    \item \textbf{問責：}當 AI 評分有誤時，誰負責？
\end{itemize}

%============================================================
\section{Cohen's Kappa 補充說明}
%============================================================

\subsection{什麼是 Cohen's Kappa？}

Cohen's Kappa ($\kappa$) 是衡量兩位評分者一致性的統計指標，校正了隨機一致的機率。

\textbf{公式：}
$$\kappa = \frac{P_o - P_e}{1 - P_e}$$

其中：
\begin{itemize}[leftmargin=2em]
    \item $P_o$ = 觀察到的一致比例 (observed agreement)
    \item $P_e$ = 隨機預期的一致比例 (expected agreement by chance)
\end{itemize}

\subsection{Quadratic-weighted Kappa}

\begin{itemize}[leftmargin=2em]
    \item 適用於有順序性的評分量表（如 1-5 分）
    \item 評分差距越大，懲罰越重（二次方權重）
    \item 比 linear-weighted kappa 對大差距更敏感
    \item 適合評估「差多少」比「對不對」更重要的情境
\end{itemize}

%============================================================
\section{參考文獻}
%============================================================

\begin{enumerate}
    \item Kuling G, Pullman S, Vasilev D, et al. Assessment of Short-Answer Questions by ChatGPT in a Medical School Course. \href{https://doi.org/10.1056/AIcs2500239}{\textit{NEJM AI}. 2026;3(2).}

    \item Roediger HL, Karpicke JD. Test-Enhanced Learning: Taking Memory Tests Improves Long-Term Retention. \textit{Psychol Sci}. 2006;17(3):249-255.

    \item McHugh ML. Interrater Reliability: The Kappa Statistic. \textit{Biochem Med (Zagreb)}. 2012;22(3):276-282.

    \item Kung TH, Cheatham M, Medenilla A, et al. Performance of ChatGPT on USMLE: Potential for AI-Assisted Medical Education Using Large Language Models. \textit{PLOS Digit Health}. 2023;2(2):e0000198.
\end{enumerate}

\end{document}
